% Options for packages loaded elsewhere
\PassOptionsToPackage{unicode}{hyperref}
\PassOptionsToPackage{hyphens}{url}
\PassOptionsToPackage{space}{xeCJK}
%
\documentclass[
]{ctexart}
\usepackage{amsmath,amssymb}
\usepackage{iftex}
\ifPDFTeX
  \usepackage[T1]{fontenc}
  \usepackage[utf8]{inputenc}
  \usepackage{textcomp} % provide euro and other symbols
\else % if luatex or xetex
  \usepackage{unicode-math} % this also loads fontspec
  \defaultfontfeatures{Scale=MatchLowercase}
  \defaultfontfeatures[\rmfamily]{Ligatures=TeX,Scale=1}
\fi
\usepackage{lmodern}
\ifPDFTeX\else
  % xetex/luatex font selection
  \setmainfont[]{Noto Serif CJK SC}
  \ifXeTeX
    \usepackage{xeCJK}
    \setCJKmainfont[]{Noto Serif CJK SC}
          \fi
  \ifLuaTeX
    \usepackage[]{luatexja-fontspec}
    \setmainjfont[]{Noto Serif CJK SC}
  \fi
\fi
% Use upquote if available, for straight quotes in verbatim environments
\IfFileExists{upquote.sty}{\usepackage{upquote}}{}
\IfFileExists{microtype.sty}{% use microtype if available
  \usepackage[]{microtype}
  \UseMicrotypeSet[protrusion]{basicmath} % disable protrusion for tt fonts
}{}
\makeatletter
\@ifundefined{KOMAClassName}{% if non-KOMA class
  \IfFileExists{parskip.sty}{%
    \usepackage{parskip}
  }{% else
    \setlength{\parindent}{0pt}
    \setlength{\parskip}{6pt plus 2pt minus 1pt}}
}{% if KOMA class
  \KOMAoptions{parskip=half}}
\makeatother
\usepackage{xcolor}
\usepackage[margin=1in]{geometry}
\usepackage{enumitem}
\setlist{itemsep=2pt, topsep=4pt}
\usepackage{titlesec}
\titlespacing*{\section}{0pt}{1.5em}{0.8em}
\titlespacing*{\subsection}{0pt}{1.2em}{0.6em}
\titlespacing*{\subsubsection}{0pt}{1.0em}{0.5em}
\usepackage{longtable,booktabs,array}
\usepackage{calc} % for calculating minipage widths
% Correct order of tables after \paragraph or \subparagraph
\usepackage{etoolbox}
\makeatletter
\patchcmd\longtable{\par}{\if@noskipsec\mbox{}\fi\par}{}{}
\makeatother
% Allow footnotes in longtable head/foot
\IfFileExists{footnotehyper.sty}{\usepackage{footnotehyper}}{\usepackage{footnote}}
\makesavenoteenv{longtable}
\setlength{\emergencystretch}{3em} % prevent overfull lines
\providecommand{\tightlist}{%
  \setlength{\itemsep}{0pt}\setlength{\parskip}{0pt}}
\setcounter{secnumdepth}{-\maxdimen} % remove section numbering
\ifLuaTeX
  \usepackage{selnolig}  % disable illegal ligatures
\fi
\usepackage{bookmark}
\IfFileExists{xurl.sty}{\usepackage{xurl}}{} % add URL line breaks if available
\urlstyle{same}
\hypersetup{
  pdftitle={DDA3020 Notes},
  pdfauthor={Hongyi Chen},
  hidelinks,
  pdfcreator={LaTeX via pandoc}}

\newcommand{\chaptercover}[3]{%
  \clearpage
  \phantomsection
  \addcontentsline{toc}{section}{#1}
  \thispagestyle{empty}
  \vspace*{0.24\textheight}
  \begin{center}
    {\LARGE\bfseries #1\par}
    \vspace{1.2em}
    {\large #2\par}
  \end{center}
  \vfill
  \noindent\textbf{Included lectures}\par
  \vspace{0.5em}
  #3
  \clearpage
}

\title{DDA3020 Notes\\
\large Course notes for CUHK-Shenzhen, Spring 2026}
\author{Hongyi Chen}
\date{}

\begin{document}
\maketitle


\chaptercover{Chapter 1: Introduction to ML}{}{%
\begin{itemize}
\item Lecture 1: Course Overview
\item Lecture 2: Probability \& Information Theory
\item Lecture 3: Linear Algebra Review
\item Lecture 4: Basic Optimization Review
\end{itemize}}

\subsection{Lecture 1: Course
Overview}\label{lecture-1-course-overview}

Learning paradigms of ML: - \textbf{Supervised Learning:} Learn from
teacher - \textbf{Unsupervised Learning:} Learn from oneself -
\textbf{Reinforcement Learning:} Learn from interaction with environment
(reward/punishment)

\begin{longtable}[]{@{}lll@{}}
\toprule\noalign{}
& \textbf{Supervised Learning} & \textbf{Unsupervised Learning} \\
\midrule\noalign{}
\endhead
\bottomrule\noalign{}
\endlastfoot
\textbf{Discrete} & Classification & Clustering \\
\textbf{Continuous} & Regression & Dimensionality Reduction \\
\end{longtable}


\clearpage

\subsection{Lecture 2: Probability \& Information
Theory}\label{lecture-2-probability-information-theory}

\begin{enumerate}
\def\labelenumi{\arabic{enumi}.}
\item
  Probability of discrete r.v.
\item
  Probability of continuous r.v.
\item
  Some popular distributions

  3.1. \textbf{Bernoulli Distribution}

  \(x \in \{0,1\}\) \(\rightarrow\) toss a coin, x = 1 (head) with
  probability \(\mu\), x = 0 (tail) with probability \(1-\mu\) .

  \[Bern(x|\mu) = \mu^x(1-\mu)^{1-x}\]

  \[E[x] = \mu\]\\
  \[Var[x] = \mu(1-\mu)\]

  3.2. \textbf{Binomial Distribution}

  Toss \(N\) coins, each follows Bernoulli distribution \(p(x|\mu)\),
  count number of heads \(m\).

  \[Bin(m|N,\mu) = \binom{N}{m} \mu^m(1-\mu)^{N-m}\]

  \[E[m] = N\mu\]\\
  \[Var[m] = N\mu(1-\mu)\]

  3.3. \textbf{Gaussian/Normal Distribution}

  \[N(x|\mu,\sigma^2) = \frac{1}{\sqrt{2\pi\sigma^2}}exp(-\frac{(x-\mu)^2}{2\sigma^2})\]\\
  For a D-dimensional vector \(\mathbf{x}\), the multivariate Gaussian
  distribution is defined as:

  \[N(\mathbf{x}|\boldsymbol{\mu},\boldsymbol{\Sigma}) = \frac{1}{(2\pi)^{D/2}|\boldsymbol{\Sigma}|^{1/2}}exp(-\frac{1}{2}(\mathbf{x}-\boldsymbol{\mu})^T\boldsymbol{\Sigma}^{-1}(\mathbf{x}-\boldsymbol{\mu}))\]
  where \(\boldsymbol{\mu}\) is the mean vector and
  \(\boldsymbol{\Sigma}\) is the covariance matrix.
\item
  Information Theory

  4.1. \textbf{Definition}\\
  For a discrete r.v. with a finite alphabet, as follows:
  \[X = \{x_1, x_2, ..., x_n\}\] We difine the \emph{Information} as
  \[I(x) = -log p(x)\]

  Information measures the amount of ``surprise'' associated with an
  event. The less likely an event is, the more information it provides
  when it occurs.

  4.2. \textbf{Entropy}\\
  Entropy is the expected value of the information content of a random
  variable. It quantifies the uncertainty associated with a random
  variable. \[H_p(X) = E[I(x_k)] = -\sum_{x_k \in X} p(x_k) log p(x_k)\]

  Apparently, entropy is maximized when all outcomes are equally likely,
  i.e., \(p(x_k) = \frac{1}{n}\) for all \(k\).

  4.3. \textbf{Cross Entropy}\\
  Cross entropy is: when the real distribution is \(P\), but we have to
  use another distribution \(Q\) to encode the information, how many
  bits are needed on average?
  \[H_{p,q}(X) = -\sum_{x_k \in X} P(X = x_k) \cdot log (Q(X = x_k))\]

  There are two properties of cross entropy:

  \begin{itemize}
  \tightlist
  \item
    Non-negativity: \(H_{P,Q}(X) \geq 0\)
  \item
    Cross entropy is not smaller than entropy:
    \(H_{P,Q}(X) \geq H_P(X)\), and the equality holds iff \(P = Q\).
  \end{itemize}

  4.4. \textbf{Relative entropy: Kullback-Leibler (KL) Divergence} KL
  divergence measures the distance between two probability
  distributions.
  \[ D_{P,Q}(X) = \sum_{x_k \in X} P(X = x_k) log \frac{P(X = x_k)}{Q(X = x_k)} = H_{P,Q}(X) - H_P(X)\]
  \[ D_{P,Q}(X) = \int_{x \in X} p_X(x) log \left(\frac{p_X(x)}{q_X(x)}\right) dx \]
  There are two properties of KL divergence:

  \begin{itemize}
  \tightlist
  \item
    Non-negativity: \(D_{P,Q}(X) \geq 0\)
  \item
    Asymmetry: \(D_{P,Q}(X) \neq D_{Q,P}(X)\)
  \end{itemize}

  The cross entropy \(H_{P,Q}(X)\) is the entropy of the distribution
  \(H_P(X)\) plus the KL divergence \(D_{P,Q}(X)\) between the two
  distributions. \[H_{P,Q}(X) = H_P(X) + D_{P,Q}(X)\]
\end{enumerate}


\clearpage

\subsection{Lecture 3: Linear Algebra
Review}\label{lecture-3-linear-algebra-review}

\begin{enumerate}
\def\labelenumi{\arabic{enumi}.}
\item
  Verctor, matrix, and their norms 1.1. Vector, matrix and their norms

  \begin{itemize}
  \tightlist
  \item
    Trace: \(tr(A) = \sum_{i} A_{ii}\)
  \item
    Vector norms:

    \begin{itemize}
    \tightlist
    \item
      \(l_0\) norm: \(\|x\|_0 = \sum_{i} I(x_i \neq 0)\), number of
      nonzero elements in vector \(x\).
    \item
      \(l_1\) norm: \(\|x\|_1 = \sum_{i} |x_i|\)
    \item
      \(l_2\) norm: \(\|x\|_2 = \sqrt{\sum_{i} x_i^2}\)
    \item
      \(l_p\) norm: \(\|x\|_p = (\sum_{i} |x_i|^p)^{1/p}\)
    \end{itemize}
  \item
    Matrix norms:

    \begin{itemize}
    \tightlist
    \item
      Frobenius norm:
      \(\|A\|_F = \sqrt{\sum_{i=1}^{m} \sum_{j=1}^{n} x_{ij}^2} = \sqrt{tr(A^TA)}\)
    \item
      Spectral norm: \(\|A\|_2 = \sigma_{max}(A)\), the largest singular
      value of \(A\). (Singular Value Decomposition:
      \(X = U\Sigma V^T\), \(\Sigma = diag(\sigma_1, \sigma_2, ...)\))
    \end{itemize}
  \end{itemize}
\item
  Matrix inverse, determinant, independence 2.1. Matrix inverse For a
  square matrix \(A\), if there exists a matrix \(B\) such that
  \(AB = BA = I\), then \(B\) is called the inverse of \(A\), denoted as
  \(A^{-1}\).

  2.2. Determinant Via SVD: \[det(A) = \prod_{i} \sigma_i\]

  2.3. Linear Independence \[c_1v_1 + c_2v_2 + ... + c_nv_n = 0\] only
  holds for \(c_1 = ... = c_n = 0\).
\item
  Systems of linear equations \[\mathbf{Xw} = \mathbf{y},\] where \[
   \mathbf{X} = \begin{bmatrix}
   x_{1,1} & \cdots & x_{1,d} \\
   \vdots & \ddots & \vdots \\
   x_{m,1} & \cdots & x_{m,d}
   \end{bmatrix}, \quad
   \mathbf{w} = \begin{bmatrix}
   w_1 \\
   \vdots \\
   w_d
   \end{bmatrix}, \quad
   \mathbf{y} = \begin{bmatrix}
   y_1 \\
   \vdots \\
   y_m
   \end{bmatrix}.
   \]
\end{enumerate}


\clearpage

\subsection{Lecture 4: Basic Optimization
Review}\label{lecture-4-basic-optimization-review}

\begin{enumerate}
\def\labelenumi{\arabic{enumi}.}
\item
  Convex set 1.1. Affine line
  \[ x = \theta x_1 + (1-\theta)x_2, \theta \in R \] 1.2. Convex set
  \[ x = \theta x_1 + (1-\theta)x_2, \theta \in [0,1] \] The convex set
  contains all line segments between any two points in the set.
\item
  Convex function 2.1. Definition \(f\): \(R^n \rightarrow R\) is convex
  if \(dom f\) is a convex set and for all \(x_1, x_2 \in dom f\) and
  \[ f(\theta x_1 + (1-\theta)x_2) \leq \theta f(x_1) + (1-\theta)f(x_2), \theta \in [0,1] \]
  2.2. First-order condition If \(f\) is differentiable, then \(f\) is
  convex iff for all \(x_1, x_2 \in dom f\),
  \[ f(y) \geq f(x) + \nabla f(x)^T (y - x) \] 2.3. Second-order
  condition If \(f\) is twice differentiable, then \(f\) is convex iff
  for all \(x \in dom f\), \[ \nabla^2 f(x) \succeq 0 \] i.e., the
  Hessian matrix is positive semidefinite. 2.4. Jensen's Inequality If f
  is convex, then for \(0 \leq \theta \leq 1\),
  \[ f(\theta x_1 + (1-\theta)x_2) \leq \theta f(x_1) + (1-\theta)f(x_2) \]
  More generally, for random variable \(X\), \[ f(E[X]) \leq E[f(X)] \]
\item
  Convex Optimization Problem A convex optimization problem has the
  form: \begin{align*}
   \text{minimize} \quad & f_0(x) \\
   \text{subject to} \quad & f_i(x) \leq 0, \quad i = 1, ..., m \\
   & h_j(x) = 0, \quad j = 1, ..., p
\end{align*} where \(f_0, f_1, ..., f_m\) are convex functions and
  \(h_1, ..., h_p\) are affine functions. 3.1. Local and Global Minimum
  \textbf{Theorem:} In a convex optimization problem, any local minimum
  is also a global minimum.
\item
  Unconstrained Minimization 4.1. Gradient Descent
  \texttt{Given\ a\ starting\ point\ x.\ \ \ \ \ repeat\ \ \ \ \ \ \ \ \ 1.\ Determine\ a\ descent\ direction\ delta\_x\ \ \ \ \ \ \ \ \ 2.\ Choose\ a\ step\ size\ t\ \textgreater{}\ 0\ \ \ \ \ \ \ \ \ 3.\ Update\ x\ :=\ x\ +\ t\ *\ delta\_x\ \ \ \ \ \ \ \ \ until\ stopping\ criterion\ is\ satisfied.}

  4.2. Line Search Method \textbf{Exact line search:}
  \(t = argmin_{t \geq 0} f(x + t \Delta x)\) \textbf{Backtracking line
  search (inexact)} Starting at \(t = 1\) Repeat \(t := \beta * t\)
  until
  \[f(x + t * \Delta x) \leq f(x) + \alpha * t * \nabla f(x)^T * \Delta x\]
\item
  Constrained Minimization: Lagrangian duality, KKT conditions 5.1.
  Lagrangian and Lagrange Dual Function Given a general constrained
  optimization problem: \begin{align*}
   \text{minimize} \quad & f_0(x) \\
   \text{subject to} \quad & f_i(x) \leq 0, \quad i = 1, ..., m \\
   & h_j(x) = 0, \quad j = 1, ..., p
\end{align*} 5.1. Lagrangian and Lagrange Dual Function The Lagrangian is
  defined as:
  \[ L(x, u, v) = f_0(x) + \sum_{i=1}^{m} u_i f_i(x) + \sum_{j=1}^{p} v_j h_j(x) \]
  The Lagrange dual function is defined as:
  \[ g(u, v) = inf_{x} L(x, u, v) \] The dual problem is: \begin{align*}
   \text{maximize} \quad & g(u, v) \\
   \text{subject to} \quad & u \succeq 0
\end{align*} 5.2. KKT Conditions A necessary condition for \(x^*\) to be
  optimal for general constrained optmization problem, a sufficient
  condition for \(x^*\) to be optimal for convex optimization problem:

  \begin{itemize}
  \tightlist
  \item
    Stationarity:
    \(0 \in \partial f(x) + \sum_{i=1}^{m} u_i \partial f_i(x) + \sum_{j=1}^{p} v_j \partial \ell_j(x)\)
  \item
    Complementary slackness: \(u_i f_i(x) = 0 \text{ for all } i\)
  \item
    Primal feasibility:
    \(h_i(x) \leq 0; \ell_j(x) = 0 \text{ for all } i,j\)
  \item
    Dual feasibility: \(u_i \geq 0 \text{ for all } i\)
  \end{itemize}
\end{enumerate}


\chaptercover{Chapter 2: Classical Supervised Learning}{Contents: Linear Regression, Logistic Regression, SVM, Decision Tree, Random Forest etc.}{%
\begin{itemize}
\item Lecture 5: Linear Regression
\item Lecture 6: Logistic Regression
\item Lecture 7: Support Vector Machine (SVM)
\item Lecture 8: Tree-based Methods
\end{itemize}}

\subsection{Lecture 5: Linear
Regression}\label{lecture-5-linear-regression}

\begin{enumerate}
\def\labelenumi{\arabic{enumi}.}
\item
  Basic Formulation Goal: Find a linear function that best fits the data
  points. \[y = \mathbf{w}^T \mathbf{x} + b\] Obejective function (MSE):
  \[\frac{1}{m} \sum_{i=1}^{m} (y_i - f_\mathbf{w}(\mathbf{x}_i))^2\]
\item
  Probabilistic Perspective Assume
  \(y = \mathbf{w}^T \mathbf{x} + b + \epsilon\), where
  \(\epsilon \sim N(0, \sigma^2)\). Thus, the output \(y\) can be seen
  as a random variable,
  \(p(y|\mathbf{x}, \mathbf{w}) = N(y|\mathbf{w}^T \mathbf{x} + b, \sigma^2)\).

  \textbf{Maximum Likelihood Estimation (MLE):}
  \[\mathbf{w}_{MLE} = \arg \max_{\mathbf{w}} \log \mathcal{L}(\mathbf{w}; D),\]

  \begin{align*}
   \log \mathcal{L}(\mathbf{w}; D) &= \log \left( \prod_{i=1}^{m} p(y_i|\mathbf{x}_i, \mathbf{w}) \right) = \sum_{i=1}^{m} \log \mathcal{N}(\mathbf{w}^\top \mathbf{x}_i, \sigma^2) \\
   &= -m \log(\sigma(2\pi)^{\frac{1}{2}}) - \frac{1}{2\sigma^2} \sum_{i=1}^{m} (y_i - \mathbf{w}^\top \mathbf{x}_i)^2.
\end{align*}

  Removing constants w.r.t. \(\mathbf{w}\), we have
  \[\mathbf{w}_{MLE} = \arg \min_{\mathbf{w}} \ \sum_{i=1}^{m} (y_i - \mathbf{w}^\top \mathbf{x}_i)^2.\]
\item
  Analytical Solution The squared error function can be written
  compactly using \(\mathbf{e} = \mathbf{Xw} - \mathbf{y}\) as: \[
   \begin{aligned}
   J(\mathbf{w}) &= \mathbf{e}^\top \mathbf{e} = (\mathbf{Xw} - \mathbf{y})^\top (\mathbf{Xw} - \mathbf{y}) \\
   &= \mathbf{w}^\top \mathbf{X}^\top \mathbf{Xw} - 2\mathbf{w}^\top \mathbf{X}^\top \mathbf{y} + \mathbf{y}^\top \mathbf{y}.
   \end{aligned}
   \] Setting the gradient to zero, we have
  \[\nabla_{\mathbf{w}} J(\mathbf{w}) = 2\mathbf{X}^\top \mathbf{Xw} - 2\mathbf{X}^\top \mathbf{y} = 0.\]
  Thus, if \(\mathbf{X}^\top \mathbf{X}\) is invertible, the closed-form
  solution is
  \[\text{Learning:} \quad \hat{\mathbf{w}} = (\mathbf{X}^\top \mathbf{X})^{-1} \mathbf{X}^\top \mathbf{y}.\]
  \[\text{Prediction:} \quad \hat{y}(X_\text{new}) = \mathbf{X}_{\text{new}}\hat{\mathbf{w}}\]
\item
  Design Matrix with Bias Term \[X = \begin{bmatrix}
   1 & x_{1,1} & x_{1,2} & \cdots & x_{1,d} \\
   1 & x_{2,1} & x_{2,2} & \cdots & x_{2,d} \\
   \vdots & \vdots & \vdots & \ddots & \vdots \\
   1 & x_{m,1} & x_{m,2} & \cdots & x_{m,d}
   \end{bmatrix}\] \[\mathbf{w} = \begin{bmatrix}
   w_0 \\
   w_1 \\
   w_2 \\
   \vdots \\
   w_d
   \end{bmatrix}\]
  \[\mathbf{x}_i^T \mathbf{w} = [1,x_1,x_2,...,x_d]_i \cdot [w_0,w_1,w_2,...,w_d]^T\]
  So, \(w_0\) is the bias term.
\item
  Gradient Descent Method Update rule:
  \[\mathbf{w} := \mathbf{w} - \eta \nabla_{\mathbf{w}} J(\mathbf{w})\]
  where \(\eta\) is the learning rate / step size, and
  \[\nabla_{\mathbf{w}} J(\mathbf{w}) = \mathbf{X}^\top (\mathbf{Xw} - \mathbf{y}).\]
\item
  Closed-form vs.~GD Closed-form: No hyperparameter or iteration, but
  needs matrix inversion \(O(d^3+m^3)\), not suitable for large-scale
  data. GD: \(O(Tmd)\), works well when d is large.
\item
  Linear Regression with multiple outputs
  \[\mathbf{f_W(X)} = \mathbf{XW}\] \[
  \begin{aligned}
  \begin{matrix}
  \text{Sample 1} \ \dashrightarrow \\
  \vdots \\
  \text{Sample m} \ \dashrightarrow
  \end{matrix}
  \underbrace{\begin{bmatrix}
  \mathbf{x}_1^T \\
  \vdots \\
  \mathbf{x}_m^T
  \end{bmatrix}}_{\mathbf{X}} \mathbf{W} 
  &= \underbrace{\begin{bmatrix}
  1 & x_{1,1} & \cdots & x_{1,d} \\
  \vdots & \vdots & \ddots & \vdots \\
  1 & x_{m,1} & \cdots & x_{m,d}
  \end{bmatrix}}_{m \times (d+1)}
  \underbrace{\begin{bmatrix}
  w_{0,1} & \cdots & w_{0,h} \\
  w_{1,1} & \cdots & w_{1,h} \\
  \vdots & \ddots & \vdots \\
  w_{d,1} & \cdots & w_{d,h}
  \end{bmatrix}}_{(d+1) \times h} \\
  \begin{matrix}
  \text{Sample 1's output} \ \dashrightarrow \\
  \vdots \\
  \text{Sample m's output} \ \dashrightarrow
  \end{matrix}
  &= \underbrace{\begin{bmatrix}
  f_{1,1} & \cdots & f_{1,h} \\
  \vdots & \ddots & \vdots \\
  f_{m,1} & \cdots & f_{m,h}
  \end{bmatrix}}_{m \times h}
  \end{aligned}
  \]

  \[\mathbf{J(W)} = \text{trace}(\mathbf(E^T E)) = \text{trace}((\mathbf{XW} - \mathbf{Y})^T(\mathbf{XW} - \mathbf{Y}))\]

  In general, the closed-form solution is
  \[\text{Learning:} \quad \hat{\mathbf{W}} = (\mathbf{X}^\top \mathbf{X})^{-1} \mathbf{X}^\top \mathbf{Y}.\]
  \[\text{Prediction:} \quad \hat{\mathbf{Y}}(\mathbf{X}_{\text{new}}) = \mathbf{X}_{\text{new}}\hat{\mathbf{W}}\]
\item
  Linear Regression for Classification Data: \({x_i,y_i}_{i=1}^m\), with
  \(x_i \in \mathcal{X}\), \(y_i \in \mathcal{Y}\). If \(\mathcal{Y}\)
  is a finiteand discrete set, then the task corresponds to
  classification.

  \begin{itemize}
  \item
    \textbf{Binary Classification} If \(\mathbf{X^TX}\) is invertible,
    the closed-form solution is
    \[\text{Learning:} \quad \hat{\mathbf{w}} = (\mathbf{X}^\top \mathbf{X})^{-1} \mathbf{X}^\top \mathbf{y}. \\
     \text{Prediction:} \quad \hat{y}(X_\text{new}) = \operatorname{sgn}(\mathbf{X}_{\text{new}}\hat{\mathbf{w}}).\]
  \item
    \textbf{Multi-class Classification}
    \[\text{Learning:} \quad \hat{\mathbf{W}} = (\mathbf{X}^\top \mathbf{X})^{-1} \mathbf{X}^\top \mathbf{Y}. \\
     \text{Prediction:} \quad \hat{\mathbf{Y}}(\mathbf{X}_{\text{new}}) = \arg\max_{i=1,...,c} \mathbf{X}_{\text{new}}\hat{\mathbf{W}}.\]
    Each row in Y is a one-hot vector. For example, if there are 3
    classes and the label is class 2, then the one-hot vector is
    {[}0,1,0{]}.
  \end{itemize}
\item
  Variants of Linear Regression 9.1. Ridge Regression (L2
  Regularization) Motivation: We can not gaurantee that
  \(\mathbf{X^TX}\) is invertible.
  \[J(\mathbf{w}) = \sum_{i=1}^{m} (y_i - \mathbf{w}^\top \mathbf{x}_i)^2 + \lambda \|\mathbf{w}\|_2^2\]
  Closed-form solution:
  \[\hat{\mathbf{w}} = (\mathbf{X}^\top \mathbf{X} + \lambda \mathbf{I})^{-1} \mathbf{X}^\top \mathbf{y}.\]
  \((X^TX + \lambda I)\) is always invertible when \(\lambda > 0\).

  9.2.Polynomial Regression Sometimes, the relationship between x and y
  is not linear. \[
   f(\mathbf{x}) = w_0 + \sum_{j=1}^{d} w_j x_j + \sum_{j=1}^{d} \sum_{k=j}^{d} w_{jk} x_j x_k + \cdots
   \]

  See an example: Data \(\{ \mathbf{x}_i, y_i \}_{i=1}^m\):

  \begin{itemize}
  \tightlist
  \item
    \(\{x_1 = 0, x_2 = 0\} \rightarrow \{y = -1\}\)
  \item
    \(\{x_1 = 1, x_2 = 1\} \rightarrow \{y = -1\}\)
  \item
    \(\{x_1 = 1, x_2 = 0\} \rightarrow \{y = +1\}\)
  \item
    \(\{x_1 = 0, x_2 = 1\} \rightarrow \{y = +1\}\)
  \end{itemize}

  \(2^{\text{nd}}\) order polynomial model: \[
   \mathbf{P} = \begin{bmatrix} 
   1 & x_1 & x_2 & x_1x_2 & x_1^2 & x_2^2 \\
   1 & x_1 & x_2 & x_1x_2 & x_1^2 & x_2^2 \\
   1 & x_1 & x_2 & x_1x_2 & x_1^2 & x_2^2 \\
   1 & x_1 & x_2 & x_1x_2 & x_1^2 & x_2^2 
   \end{bmatrix} =
   \begin{bmatrix} 
   1 & 0 & 0 & 0 & 0 & 0 \\
   1 & 1 & 1 & 1 & 1 & 1 \\
   1 & 1 & 0 & 0 & 1 & 0 \\
   1 & 0 & 1 & 0 & 0 & 1 
   \end{bmatrix}
   \]

  \(\hat{\mathbf{w}} = \mathbf{P}^\top (\mathbf{P}\mathbf{P}^\top)^{-1} \mathbf{y}\)
  (note: under determined linear system)

  \[
   = \begin{bmatrix}
   1 & 1 & 1 & 1 \\
   0 & 1 & 1 & 0 \\
   0 & 1 & 0 & 1 \\
   0 & 1 & 0 & 0 \\
   0 & 1 & 1 & 0 \\
   0 & 1 & 0 & 1
   \end{bmatrix}
   \begin{bmatrix}
   1 & 1 & 1 & 1 \\
   1 & 6 & 3 & 3 \\
   1 & 3 & 3 & 1 \\
   1 & 3 & 1 & 3
   \end{bmatrix}^{-1}
   \begin{bmatrix}
   -1 \\ -1 \\ +1 \\ +1
   \end{bmatrix} =
   \begin{bmatrix}
   -1 \\ 1 \\ 1 \\ -4 \\ 1 \\ 1
   \end{bmatrix}
   \]

  9.3. Lasso Regression (L1 Regularization)
  \[J(\mathbf{w}) = \sum_{i=1}^{m} (y_i - \mathbf{w}^\top \mathbf{x}_i)^2 + \lambda \|\mathbf{w}\|_1\]

  9.4. Robust Regression
  \[J(\mathbf{w}) = \sum_{i=1}^{m} |y_i - \mathbf{w}^\top \mathbf{x}_i|\]
\end{enumerate}

\textbf{Summary of Probabilistic Perspectives:}

\begin{longtable}[]{@{}
  >{\raggedright\arraybackslash}p{(\columnwidth - 4\tabcolsep) * \real{0.3333}}
  >{\raggedright\arraybackslash}p{(\columnwidth - 4\tabcolsep) * \real{0.3333}}
  >{\raggedright\arraybackslash}p{(\columnwidth - 4\tabcolsep) * \real{0.3333}}@{}}
\toprule\noalign{}
\begin{minipage}[b]{\linewidth}\raggedright
\(p(y \mid \mathbf{x}, \mathbf{w})\)
\end{minipage} & \begin{minipage}[b]{\linewidth}\raggedright
\(p(\mathbf{w})\)
\end{minipage} & \begin{minipage}[b]{\linewidth}\raggedright
regression method
\end{minipage} \\
\midrule\noalign{}
\endhead
\bottomrule\noalign{}
\endlastfoot
Gaussian & Uniform & Least squares \\
Gaussian & Gaussian & Ridge regression \\
Gaussian & Laplace & Lasso regression \\
Laplace & Uniform & Robust regression \\
Student & Uniform & Robust regression \\
\end{longtable}


\clearpage

\subsection{Lecture 6: Logistic
Regression}\label{lecture-6-logistic-regression}

\paragraph{1. Classification}\label{classification}

Sigmoid function: \[\sigma(z) = \frac{1}{1 + e^{-z}}\]

\(f_\mathbf{w}(\mathbf{x})\) is the probability of \(y=1\) given
\(\mathbf{x}\) and \(\mathbf{w}\), which is defined as:
\[p(y=1|\mathbf{x}, \mathbf{w}) = \sigma(\mathbf{w}^\top \mathbf{x})\]

\paragraph{2. Logistic Regression}\label{logistic-regression}

\textbf{Linear vs.~Logistic Regression:}

\begin{longtable}[]{@{}
  >{\raggedright\arraybackslash}p{(\columnwidth - 4\tabcolsep) * \real{0.3333}}
  >{\raggedright\arraybackslash}p{(\columnwidth - 4\tabcolsep) * \real{0.3333}}
  >{\raggedright\arraybackslash}p{(\columnwidth - 4\tabcolsep) * \real{0.3333}}@{}}
\toprule\noalign{}
\begin{minipage}[b]{\linewidth}\raggedright
Aspect
\end{minipage} & \begin{minipage}[b]{\linewidth}\raggedright
Linear Regression
\end{minipage} & \begin{minipage}[b]{\linewidth}\raggedright
Logistic Regression
\end{minipage} \\
\midrule\noalign{}
\endhead
\bottomrule\noalign{}
\endlastfoot
Loss Function & MSE (Mean Squared Error) & Cross-Entropy Loss \\
Formula &
\(\frac{1}{2m} \sum_{i=1}^{m}(f_\mathbf{w}(\mathbf{x}^{(i)}) - y^{(i)})^2\)
&
\(-y \log(f_\mathbf{w}(\mathbf{x})) - (1-y) \log(1-f_\mathbf{w}(\mathbf{x}))\) \\
Output Range & \((-\infty, +\infty)\) & \([0, 1]\) (probability) \\
Convexity & Convex & Non-convex with MSE; Convex with Cross-Entropy \\
\end{longtable}

\textbf{Why Cross-Entropy for Logistic Regression?}

\begin{enumerate}
\def\labelenumi{\arabic{enumi}.}
\item
  \textbf{MSE Problems:}

  \begin{itemize}
  \tightlist
  \item
    Non-convexity
  \item
    Gradient vanishing
  \end{itemize}
\item
  \textbf{Cross-Entropy Advantages:}
  \[\text{cost} = \begin{cases} -\log(f_\mathbf{w}(\mathbf{x})) & \text{if } y=1 \\ -\log(1-f_\mathbf{w}(\mathbf{x})) & \text{if } y=0 \end{cases}\]

  \begin{itemize}
  \tightlist
  \item
    \textbf{Rewards correct predictions:} Loss approaches 0 when
    \(f_\mathbf{w}(\mathbf{x})\) matches \(y\).
  \item
    \textbf{Penalizes errors:} Loss approaches \(\infty\) when
    predictions are far from true labels, forcing rapid parameter
    correction.
  \end{itemize}
\end{enumerate}

\paragraph{3. Gradient Descent for Logistic
Regression}\label{gradient-descent-for-logistic-regression}

Learning \(\mathbf{w}\) by minimizing the cost function:
\[\mathbf{w}^* = \arg \min_{\mathbf{w}} J(\mathbf{w}) = \arg \min_{\mathbf{w}} -\frac{1}{m} \sum_{i=1}^{m} [y_i \log(f_\mathbf{w}(\mathbf{x}_i)) + (1-y_i) \log(1-f_\mathbf{w}(\mathbf{x}_i))]\]

Repeat the following steps until convergence:
\[\mathbf{w} := \mathbf{w} - \eta \nabla_{\mathbf{w}} J(\mathbf{w})\]

\[\nabla_{\mathbf{w}} J(\mathbf{w}) = \frac{1}{m} \sum_{i=1}^{m} (f_\mathbf{w}(\mathbf{x}_i) - y_i) \mathbf{x}_i\]

\paragraph{4. Multiclass Logistic Regression (Softmax
Regression)}\label{multiclass-logistic-regression-softmax-regression}

\textbf{Implementation Strategies}

\begin{longtable}[]{@{}
  >{\raggedright\arraybackslash}p{(\columnwidth - 4\tabcolsep) * \real{0.3333}}
  >{\raggedright\arraybackslash}p{(\columnwidth - 4\tabcolsep) * \real{0.3333}}
  >{\raggedright\arraybackslash}p{(\columnwidth - 4\tabcolsep) * \real{0.3333}}@{}}
\toprule\noalign{}
\begin{minipage}[b]{\linewidth}\raggedright
Strategy
\end{minipage} & \begin{minipage}[b]{\linewidth}\raggedright
One-vs-All (OvR)
\end{minipage} & \begin{minipage}[b]{\linewidth}\raggedright
Softmax Regression (Multinomial)
\end{minipage} \\
\midrule\noalign{}
\endhead
\bottomrule\noalign{}
\endlastfoot
\textbf{Core Idea} & Train \(C\) binary classifiers (one per class). &
Train one unified model for \(C\) classes. \\
\textbf{Logic} & Predict class with highest binary probability. & Output
a normalized probability distribution. \\
\textbf{Use Case} & Overlapping categories (e.g., multi-tagging). &
Mutually exclusive categories (single label). \\
\end{longtable}

\textbf{Softmax Function} For input \(\mathbf{x}\), the probability of
belonging to class \(j\) is:
\[P(y = j \mid \mathbf{x}; \mathbf{W}) = \frac{\exp(\mathbf{w}_j^T \mathbf{x})}{\sum_{c=1}^C \exp(\mathbf{w}_c^T \mathbf{x})}\]

\textbf{Cost Function and Optimization} \textbf{Cross-Entropy Loss:}
\[J(\mathbf{W}) = -\frac{1}{m} \sum_{i=1}^m \sum_{j=1}^C [\mathbb{I}(y_i = j) \log(f_{\mathbf{w}_j}(\mathbf{x}_i))]\]
Where \(\mathbb{I}(\cdot)\) is the indicator function (effectively
one-hot encoding).

\textbf{Gradient Descent:}
\[\nabla_{\mathbf{w}_j} J(\mathbf{W}) = \frac{1}{m} \sum_{i=1}^m [f_{\mathbf{w}_j}(\mathbf{x}_i) - \mathbb{I}(y_i = j)] \mathbf{x}_i\]

\paragraph{5. Regularization in Logistic
Regression}\label{regularization-in-logistic-regression}

Add an \(L_2\) norm penalty term to the original cost function \(J(w)\):
\[\bar{J}(w) = J(w) + \frac{\lambda}{2m} \sum_{j=1}^d w_j^2\]

\begin{itemize}
\tightlist
\item
  \textbf{\(\lambda\) (Regularization parameter):} Controls the penalty
  strength. A larger \(\lambda\) makes the model simpler (prevents
  overfitting), but a value too large can cause underfitting.
\item
  \textbf{No penalty for \(w_0\):} The bias term \(w_0\) is usually not
  regularized because it only translates the decision boundary and does
  not affect the model's complexity.
\end{itemize}

\textbf{Gradient Descent Update Rule:} The parameter updates with
regularization are slightly modified:

For \(w_0\) (unchanged):
\[w_0 \leftarrow w_0 - \alpha \frac{1}{m} \sum_{i=1}^m (f_w(x_i) - y_i) x_i^{(0)}\]

For \(w_j\) (\(j > 0\)):
\[w_j \leftarrow w_j - \alpha \left[ \frac{1}{m} \sum_{i=1}^m (f_w(x_i) - y_i) x_i^{(j)} + \frac{\lambda}{m} w_j \right]\]


\clearpage

\subsection{Lecture 7: Support Vector Machine
(SVM)}\label{lecture-7-support-vector-machine-svm}

\paragraph{1. Two Types of Derivation}\label{two-types-of-derivation}

\textbf{Perspective 1: Geometric Intuition and Large Margin} -
\textbf{The Distance from a Point to a Hyperplane:} For a hyperplane
defined by \(\mathbf{w}^T \mathbf{x} + b = 0\), the distance from a
point \(\mathbf{x}\) to the hyperplane is given by:
\[\text{Distance} = \frac{|\mathbf{w}^T \mathbf{x} + b|}{\|\mathbf{w}\|}\]
- \textbf{Property of the Margin:} If we scale \(\mathbf{w}\) and \(b\)
by a positive constant \(c\), the location of the hyperplane does not
change:
\[\frac{|(c\mathbf{w})^T \mathbf{x} + cb|}{\|c\mathbf{w}\|} = \frac{|\mathbf{w}^T \mathbf{x} + b|}{\|\mathbf{w}\|}\]
- \textbf{Normalization:} We can normalize \(\mathbf{w}\) and \(b\) such
that the closest points (support vectors) satisfy:
\[y_i (\mathbf{w}^T \mathbf{x}_i + b) = 1\] - \textbf{Maximizing the
Margin:} Margin over all training data points is
\[\gamma = \min_{i} \frac{|f_{\mathbf{w},b}(\mathbf{x}_i)|}{\|\mathbf{w}\|}\]
Recall \(y_i \in \{-1, +1\}\), we have:
\[\gamma = \min_{i} \frac{y_i f_{\mathbf{w},b}(\mathbf{x}_i)}{\|\mathbf{w}\|}\]
Maximize margin:
\[\max_{\mathbf{w}, b} \gamma = \max_{\mathbf{w}, b} \min_{i} \frac{y_i f_{\mathbf{w},b}(\mathbf{x}_i)}{\|\mathbf{w}\|} = \max_{\mathbf{w}, b} \min_{i} \frac{y_i (\mathbf{w}^T \mathbf{x}_i + b)}{\|\mathbf{w}\|}\]
Recall the fixed scaleing \(y_i (\mathbf{w}^T \mathbf{x}_i + b) = 1\),
we have:
\[\max_{\mathbf{w}, b} \gamma = \max_{\mathbf{w}, b} \frac{1}{\|\mathbf{w}\|} = \min_{\mathbf{w}, b} \frac{1}{2} \|\mathbf{w}\|^2\]
Subject to the constraints:
\[y_i (\mathbf{w}^T \mathbf{x}_i + b) \geq 1, \quad i = 1, ..., m\]

\textbf{Perspective 2: Hinge Loss} When \(y_i = + 1\), if the prediction
value \(f_{\mathbf{w},b}(\mathbf{x}_i) \geq 1\), the model is correct
ans \(Loss = 0\); if \(f_{\mathbf{w},b}(\mathbf{x}_i) < 1\), there will
be a linear penalty.

\textbf{Hinge Loss:} \[\max(0, 1 - y_i f_{\mathbf{w},b}(\mathbf{x}_i))\]

\subsubsection{2. Optimizing SVM by Lagrangian
Duality}\label{optimizing-svm-by-lagrangian-duality}

The objective function of SVM is:
\[\begin{aligned}\min_{\mathbf{w}, b} \quad & \frac{1}{2} \|\mathbf{w}\|^2 \\
\text{subject to} \quad & y_i (\mathbf{w}^T \mathbf{x}_i + b) \geq 1, \forall i
\end{aligned}\]

It can be transformed to:
\[\begin{aligned}\min_{\mathbf{w}, b} \quad & \frac{1}{2} \|\mathbf{w}\|^2 \\
\text{subject to} \quad & 1 - y_i (\mathbf{w}^T \mathbf{x}_i + b) \leq 0, \forall i
\end{aligned}\] The Lagrangian is:
\[L(\mathbf{w}, b, \mathbf{u}) = \frac{1}{2} \|\mathbf{w}\|^2 + \sum_{i=1}^{m} u_i (1 - y_i (\mathbf{w}^T \mathbf{x}_i + b))\]

The primal and dual optimal solutions should satisfy the KKT conditions:
- \textbf{Stationarity:}
\[\frac{\partial L}{\partial \mathbf{w}} = 0 \Rightarrow \mathbf{w} = \sum_{i}^m \alpha_i y_i \mathbf{x}_i \\
    \frac{\partial L}{\partial b} = 0 \Rightarrow \sum_{i}^m \alpha_i y_i = 0\]
- \textbf{Feasibility:}
\[\alpha_i \geq 0, \ 1 - y_i(\mathbf{w}^\top \mathbf{x}_i + b) \leq 0, \ \forall i\]
- \textbf{Complementary slackness:}
\[\alpha_i(1 - y_i(\mathbf{w}^\top \mathbf{x}_i + b)) = 0, \ \forall i\]

Replacing the stationary condition into Lagrange function, we have: \[
\begin{aligned}
\mathcal{L}(\mathbf{w}, b, \boldsymbol{\alpha}) &= \frac{1}{2}\|\mathbf{w}\|^2 + \sum_{i}^m \alpha_i - \sum_{i}^m \alpha_i y_i \left(\sum_{j}^m \alpha_j y_j \mathbf{x}_j\right)^\top \mathbf{x}_i - \sum_{i}^m \alpha_i y_i b \\
&= \frac{1}{2}\|\mathbf{w}\|^2 + \sum_{i}^m \alpha_i - \sum_{i,j}^m \alpha_i \alpha_j y_i y_j {\mathbf{x}_i^\top \mathbf{x}_j} - b \sum_{i}^m \alpha_i y_i \\
&= \sum_{i}^m \alpha_i - \frac{1}{2}\|\mathbf{w}\|^2 \\
&= \sum_{i}^m \alpha_i - \frac{1}{2} \sum_{i,j}^m \alpha_i \alpha_j y_i y_j {\mathbf{x}_i^\top \mathbf{x}_j}
\end{aligned}
\] Thus, the dual problem is:
\[\begin{aligned}\max_{\boldsymbol{\alpha}} \quad & \sum_{i}^m \alpha_i - \frac{1}{2} \sum_{i,j}^m \alpha_i \alpha_j y_i y_j {\mathbf{x}_i^\top \mathbf{x}_j} \\
\text{s.t.} \quad & \alpha_i \geq 0, \forall i \\
& \sum_{i}^m \alpha_i y_i = 0
\end{aligned}\] \textbf{It can be solved by any off-the-shelf
optimization solver.}

We use the solved \(\alpha\) back into the stationary condition:
\[\mathbf{w} = \sum_{i}^m \alpha_i y_i \mathbf{x}_i\] For any support
vector \(\mathbf{x}_j\) (i.e., \(\alpha_j > 0, j \in \mathcal{S}\)), we
have:
\[y_j (\mathbf{w}^T \mathbf{x}_j + b) = 1 \Rightarrow y_j (\sum_{i}^m \alpha_i y_i \mathbf{x}_i^T \mathbf{x}_j + b) = 1\]

Product \(y_j\) for both sides of the above equation, and utilizing
\(y_j \cdot y_j = 1\), we have:
\[\sum_{i}^m \alpha_i y_i \mathbf{x}_i^T \mathbf{x}_j + b = y_j \\
\Rightarrow b = \frac{1}{|\mathcal{S}|} \sum_{j \in \mathcal{S}} (y_j - \sum_{i}^m \alpha_i y_i \mathbf{x}_i^T \mathbf{x}_j)\]

\textbf{Prediction:}
\[\mathbf{w}^T \mathbf{x} + b = \sum_{i}^m \alpha_i y_i \mathbf{x}_i^T \mathbf{x} + \frac{1}{|\mathcal{S}|} \sum_{j \in \mathcal{S}} (y_j - \sum_{i}^m \alpha_i y_i \mathbf{x}_i^T \mathbf{x}_j)\]

If \(\mathbf{w}^T \mathbf{x} + b > 0\), then the predicted class of
\(\mathbf{x}\) is \(+1\), otherwise \(-1\).

The prediction is correct iff \(y(\mathbf{w}^T \mathbf{x} + b) > 0\).


\clearpage

\subsection{Lecture 8: Tree-based
Methods}\label{lecture-8-tree-based-methods}

前面学过的模型大多是参数模型：训练时先选定一个整个输入空间共享的模型形式，再用全部训练数据去学一组固定参数；测试时所有样本都用同一套模型、同一套参数。

它们的问题是：

\begin{itemize}
\tightlist
\item
  模型形式由人预先假设，可能和真实关系差很远；
\item
  决策过程通常不够直观，不容易解释。
\end{itemize}

引入\textbf{非参数模型} Typical nonparametric models include k-nearest
neighbors (KNN) and decision tree (DT).

\paragraph{1. Classification Tree}\label{classification-tree}

对同一训练集，可能有很多棵零训练误差的树。

One input variable each step: how to choose? * Random * Least-values *
Most-values * Impurity Measure
(选择使不纯度下降最多，即信息增益最大的属性)

Impurity 不纯度： 可用 Entropy 衡量。
\[\text{Info}(S) = -\sum_{i=1}^c p_i \log_2 p_i\]

Information Gain 信息增益： 选择使不纯度下降最多的属性。
\[\text{Gain}(S, A) = \text{Info}(S) - \text{Info}(S|A)\]

\[\text{Info}(S|A) = \sum_{v=1}^V \frac{|D_v|}{|S|} \text{Info}(D_v)\]

也可以用Gini Index 衡量 impurity。
\[\text{Gini}(S) = 1 - \sum_{i=1}^c p_i^2\]
\[\text{Gini}(S|A) = \sum_{v=1}^V \frac{|D_v|}{|S|} \text{Gini}(D_v)\]
\[ \Delta \text{Gini}(S) = \text{Gini}(S) - \text{Gini}(S|A)\]

\paragraph{2. Regression Tree}\label{regression-tree}

回归树和分类树构造方式很像，只是分类树关心类别纯度；回归树关心数值预测误差。

叶节点预测值：\(\bar{y}_c = \frac{\sum_{i \in c} y_i}{N_c}\)
回归树通过MSE/SSE来评价效果： 某个节点的MSE：
\(e_c = \frac{1}{N_c} \sum_{i \in c} (y_i - \bar{y}_c)^2\)
整棵树一般看总SSE：
\(S = \sum_{c \in \text{leaves(Tree)}} \sum_{i \in c} (y_i - \bar{y}_c)^2\)

算法： 1. 从包含全部样本的单个节点开始； 2.
对每个节点，尝试所有变量、所有二分阈值； 3. 选择使 SSE 降低最多的划分；
4. 若降低不够大、或子节点样本过少、或特征已无差异，则停止； 5.
对新子节点继续递归。

也就是说，对于某个值\(w_1\)把根节点分成左右节点\(c_L, c_R\)，则
\[\bar{y}_L = \frac{\sum_{i \in c_L} y_i}{N_L}, \quad \bar{y}_R = \frac{\sum_{i \in c_R} y_i}{N_R} \\
S_{w_1} = \sum_{i \in c_L} (y_i - \bar{y}_L)^2 + \sum_{i \in c_R} (y_i - \bar{y}_R)^2\\
w_1^* = \arg \max_{w_1} (S-S_{w_1})\]

\paragraph{3. Pruning}\label{pruning}

过程： 1. 在训练集上先长一棵比较深的树； 2.
然后反复评估``把某个子树剪掉''是否提升验证集表现； 3.
贪心地剪掉最有利于验证集准确率的节点； 4. 直到继续剪枝会带来伤害。

\paragraph{4. Ensemble Methods}\label{ensemble-methods}

\textbf{Bagging}

第一步：Bootstrap 重采样 * 从训练集有放回抽样，得到多个不同的训练子集。

第二步：每个子集训练一棵过度生长的树

第三步：聚合预测 * 分类：多数表决 * 回归：平均值

\textbf{Random Forest} 在 Bagging 的基础上，每次分裂节点时不在全部 \(N\)
个特征里找最优划分，而是随机选择 \(m\)
个特征进行划分，增加树之间的差异性。


\chaptercover{Chapter 3: Deep Learning}{Contents: Neural Networks(MLP,CNN), RNN, Transformer etc.}{%
\begin{itemize}
\item Lecture 9: Neural Networks
\item Lecture 10: Convolutional Neural Networks
\end{itemize}}

\subsection{Lecture 9: Neural
Networks}\label{lecture-9-neural-networks}

\paragraph{1. Introduction}\label{introduction}

神经单元：\(y = \sigma(\mathbf{w}^T \mathbf{x} + b)\)

Some Activation Functions: * ReLU: \(\sigma(z) = \max(0, z)\) * Soft
ReLU: \(\sigma(z) = \log(1 + e^z)\) * Sigmoid/Logistic:
\(\sigma(z) = \frac{1}{1 + e^{-z}}\) * Tanh:
\(\sigma(z) = \frac{e^z - e^{-z}}{e^z + e^{-z}}\)

\paragraph{2. Perceptron}\label{perceptron}

能表示 AND、OR、NOT 等线性可分的函数，但不能表示 XOR
等线性不可分的函数。 Formula: \(y = \operatorname{sgn}(\mathbf{w}^T \mathbf{x} + b)\)
Objective function:
\(J(\mathbf{w}, b) = -\sum_{i=1}^m y_i (\mathbf{w}^T \mathbf{x}_i + b)\)
Gradient Descent:
\(\mathbf{w} \leftarrow \mathbf{w} + \eta (y-t) \mathbf{x}\)

\paragraph{3. Multilayer Feedforward Neural
Networks}\label{multilayer-feedforward-neural-networks}

一个多层前馈神经网络包含： - 输入层 input layer - 一个或多个隐藏层
hidden layer(s) - 输出层 output layer

它的结构特点： - 信息只沿一个方向传播：从输入到输出； -
相邻层之间全连接； - 同层神经元之间没有连接； -
非相邻层之间没有连接，也没有 skip connection。

\begin{verbatim}
x -> h -> y
\end{verbatim}

多层模型： \[h = g_2(\mathbf{W}^T \mathbf{x} + \mathbf{c}) \\
y = g_1(w^T h + b)\]

\paragraph{4. Backpropagation}\label{backpropagation}

在计算图上应用链式法则


\clearpage

\subsection{Lecture 10: Convolutional Neural
Networks}\label{lecture-10-convolutional-neural-networks}

\paragraph{1. Covolution Layer}\label{covolution-layer}

Convolution operation: \(y = w * x + b\)
每一个filter（卷积核）会产生一张activation
map（特征图），每个位置的值是卷积核和输入图像对应区域的加权和。

若输入是：32×32×3，用 6 个 5×5×3 的filters，stride=1，且无
padding，那么每个 filter 生成一张 28×28 的图。 最后堆叠得到：28×28×6

输出深度 = filter 个数。

\textbf{卷积层的空间尺寸计算公式：} \[\frac{N - F + 2P}{S} + 1\]
其中，\(N\) 是输入的空间尺寸，\(F\) 是 filter 的空间尺寸，\(P\) 是
padding 的像素数，\(S\) 是 stride。

\textbf{卷积层参数量} 设输入是：\(W_1 \times H_1 \times C\) 卷积层有 4
个超参数： * \(K\)：滤波器个数 * \(F\)：滤波器尺寸 * \(S\)：stride *
\(P\)：padding

则输出为：\(W_2 \times H_2 \times K\) 其中： \[
W_2 = \frac{W_1 - F + 2P}{S} + 1
\]

\[
H_2 = \frac{H_1 - F + 2P}{S} + 1
\]

参数量为：\(F^2 \times C \times K + K\) 即： * \(F^2 \times C \times K\)
个权重 * \(K\) 个偏置

\paragraph{2. Pooling Layer}\label{pooling-layer}

MAX pooling：取局部区域的最大值，保留纹理特征。

\paragraph{3. Fully Connected Layer}\label{fully-connected-layer}

全连接层：每个神经元与上一层的所有神经元相连，输出一个标量。

\chaptercover{Chapter 4: Unsupervised Learning}{Contents: K-means, GMM, EM, PCA etc.}{%
\begin{itemize}
\item Lecture 11: RNNs
\item Lecture 12/13: Over/Under-fitting, Bias - Variance Tradeoff and Performance Evaluation
\item Lecture 14: Intro to Unsupervised Learning
\item Lecture 15: K-means Clustering
\item Lecture 16: Gaussian Mixture Models and Expectation-Maximization Algorithm
\end{itemize}}

\subsection{Lecture 11: RNNs}\label{lecture-11-rnns}

\paragraph{1. Basic RNN}\label{basic-rnn}

RNN is a type of neural network designed for sequential data.\\
序列数据的核心特征是：顺序重要、长度可变。

\textbf{''当前时刻的表示，不仅取决于当前输入，还取决于过去的状态。}

\[h_t = f(h_{t-1}, x_t)\]

这里，\(h_t\) 是当前时刻的隐藏状态，\(h_{t-1}\)
是前一时刻的隐藏状态，\(x_t\) 是当前时刻的输入，\(f\)
是一个非线性函数（如tanh或ReLU）。

e.g.~ \[h_t = \tanh(W_{hh} h_{t-1} + W_{xh} x_t)\\
y_t = W_{hy} h_t\]

\textbf{RNN 参数共享}： - \(W_{hh}\)：隐藏状态之间的权重矩阵 -
\(W_{xh}\)：输入到隐藏状态的权重矩阵 -
\(W_{hy}\)：隐藏状态到输出的权重矩阵。

Cost function:
\[E(\theta) = \frac{1}{T} \sum_{t=1}^{T} L(y_t, \hat{y}_t)\] where
\(\theta = {W,W'}\).

\textbf{Training}

RNN 的优化仍然是``反向传播 +
梯度下降''，但由于网络跨越时间展开，因此要使用 Backpropagation Through
Time，BPTT。

实际上一般使用 Truncated BPTT （截断
BPTT）来减少计算量，即只反向传播固定数量的时间步。

问题：RNN 的梯度在时间上反向传播时，会不断乘上 \(W_{hh}\)。 - 如果
\(W_{hh}\) 的特征值小于1，梯度会逐渐变小，导致梯度消失； -
如果特征值大于1，梯度会逐渐变大，导致梯度爆炸。

-\textgreater{} RNN 难以捕捉 long-term dependencies.

\paragraph{2. LSTM (optional)}\label{lstm-optional}

LSTM（Long Short-Term Memory）解决传统 RNN 的梯度消失问题。

加入了三个门控机制： - \textbf{遗忘门}：决定保留多少过去的信息。
\[f_t = \sigma(W_f [h_{t-1}, x_t] + b_f)\] -
\textbf{输入门}：决定当前输入的重要性。
\[i_t = \sigma(W_i [h_{t-1}, x_t] + b_i)\] -
\textbf{输出门}：决定输出多少当前状态的信息。
\[o_t = \sigma(W_o [h_{t-1}, x_t] + b_o)\]

LSTM 还引入了一个新的状态变量 \(C_t\)，称为细胞状态，用于存储长期信息。
\[C_t = f_t * C_{t-1} + i_t * \tilde{C}_t\] \[h_t = o_t * \tanh(C_t)\]

\paragraph{3. Limitations of RNN}\label{limitations-of-rnn}

\begin{itemize}
\tightlist
\item
  \textbf{计算效率}：RNN
  需要逐步处理序列数据，无法并行化，导致训练时间较长。
\item
  \textbf{长距离依赖}：虽然 LSTM 和 GRU
  改善了这一问题，但在处理非常长的序列时仍然存在困难。
\end{itemize}

\paragraph{4. Transformer}\label{transformer}

Transformer 的核心是 self-attention
机制，它允许模型在处理序列时关注不同位置的信息。 - \textbf{Query, Key,
Value}：每个输入位置都会生成三个向量：
\[Q = W_Q x, K = W_K x, V = W_V x\] - \textbf{Attention Score}：计算
Query 和 Key 之间的相似度：
\[\text{Attention}(Q, K, V) = \text{softmax}\left(\frac{QK^T}{\sqrt{d_k}}\right) V\]


\clearpage

\subsection{Lecture 12/13: Over/Under-fitting, Bias - Variance
Tradeoff and Performance
Evaluation}\label{lecture-1213-overunder-fitting-bias---variance-tradeoff-and-performance-evaluation}

\paragraph{Bias-Variance Tradeoff}\label{bias-variance-tradeoff}

假设数据生成过程为
\[y = f(x) + \epsilon, \epsilon \sim \mathcal{N}(0, \sigma^2)\]

机器学习的目标是基于训练集\(D\)，通过某些学习算法\(\mathcal{A}\)
（如SVM，Logistic Regression），学习一个 hypothesis function \(h\)。即
\(h_D = \mathcal{A}(D)\)。

\textbf{Expected Hypothesis Function} (given \(\mathcal{A}\)):
\[\bar{h} = \mathbb{E}_D[h_D]\\
\mathbb{E}_{y|x}[y] = f(x)\] 即在所有可能的训练集 \(D\) 上，学习算法
\(\mathcal{A}\) 学习到的假设函数的平均。

\textbf{Expected test error}:
\[\mathbb{E}_{(x,y),D}[(h_D(x)-y)^2] = \mathbb{E}_D[(h_D(x)-\bar{h}(x))^2] + \mathbb{E} [(\bar{h}(x)-f(x))^2] + \mathbb{E}[(f(x)-y)^2]\]
即 = Variance + Bias\^{}2 + Noise

\begin{itemize}
\tightlist
\item
  模型复杂度越高，Variance 越大，Bias 越小。
\item
  模型复杂度越低，Variance 越小，Bias 越大。
\end{itemize}

\paragraph{Cross Validation}\label{cross-validation}

How to tune the hyper-parameters of a model?

K-fold Cross Validation:

\begin{verbatim}
｜fold 1 ｜fold 2 ｜fold 3 ｜fold 4 ｜fold 5 | test |
\end{verbatim}

\begin{enumerate}
\def\labelenumi{\arabic{enumi}.}
\tightlist
\item
  将训练数据划分为 K 个子集。
\item
  每次用其中 1 份做 validation，其余 K-1 份做 train.
\item
  重复K次，计算平均 validation performance.
\item
  选择平均 performance 最好的 hyper-parameters.
\end{enumerate}

\textbf{为什么比分成 train/validation/test 更好？}
因为它降低了``随机切一刀''的偶然性。

如果只做一次 train/val split，某些样本恰好被分进
validation，可能会让结果偏高或偏低。 K-fold 的思想是：

让每个样本都轮流做一次验证集样本。

这样对数据的利用更充分，也更稳定。

但成本也高。

\paragraph{Evaluation Metrics for
Regression}\label{evaluation-metrics-for-regression}

\begin{itemize}
\tightlist
\item
  \textbf{Mean Absolute Error (MAE)}: 平均绝对误差
  \[MAE = \frac{1}{n} \sum_{i=1}^{n} |y_i - \hat{y}_i|\]
\item
  \textbf{Mean Squared Error (MSE)}: 平均平方误差
  \[MSE = \frac{1}{n} \sum_{i=1}^{n} (y_i - \hat{y}_i)^2\]
\end{itemize}

\paragraph{Evaluation Metrics for
Classification}\label{evaluation-metrics-for-classification}

Confusion Matrix \textbar{} \textbar{} Predicted Positive \textbar{}
Predicted Negative \textbar{}
\textbar---------------\textbar--------------------\textbar--------------------\textbar{}
\textbar{} Actual Positive \textbar{} True Positive (TP) \textbar{}
False Negative (FN) \textbar{} \textbar{} Actual Negative \textbar{}
False Positive (FP) \textbar{} True Negative (TN) \textbar{}

\begin{itemize}
\tightlist
\item
  \textbf{Accuracy}: 准确率 (所有样本中，有多少被正确分类。)
  \[Accuracy = \frac{TP + TN}{TP + TN + FP + FN}\]
\item
  \textbf{Precision}: 精确率
  (所有被模型预测成正类的样本中，真正是正类的比例。)
  \[Precision = \frac{TP}{TP + FP}\]
\item
  \textbf{Recall}: 召回率 (所有真实正类中，被模型成功找出来的比例。)
  \[Recall = \frac{TP}{TP + FN}\]
\item
  \textbf{Cost-sensitive metrics}: 不同类型的错误有不同的成本。
  对于两种错误：FP 和 FN，分配不同的权重。
  \[Cost = \frac{C_{p,p} \cdot TP + C_{n,n} \cdot TN}{C_{p,p} \cdot TP + C_{n,n} \cdot TN + C_{p,n} \cdot FP + C_{n,p} \cdot FN}\]
\end{itemize}

可以先对 Confusion Matrix 作 Normalization： \textbar{} \textbar{}
Predicted Positive \textbar{} Predicted Negative \textbar{}
\textbar---------------\textbar--------------------\textbar--------------------\textbar{}
\textbar{} Actual Positive \textbar{} TPR \textbar{} FNR \textbar{}
\textbar{} Actual Negative \textbar{} FPR \textbar{} TNR \textbar{}

\begin{itemize}
\tightlist
\item
  (True Positive Rate) TPR = TP/(TP + FN) = Recall
\item
  (False Negative Rate) FNR = FN/(TP + FN)
\item
  (False Positive Rate) FPR = FP/(FP + TN)
\item
  (True Negative Rate) TNR = TN/(FP + TN)
\end{itemize}

每一个阈值设置，都对应一组不同的 FPR 和 FNR，这样的点称为一个 operating
point

随着阈值从 0 变到 1，FPR 呈下降趋势，FNR 呈上升趋势。 他们的交点称为
Equal Error Rate (EER)。

\paragraph{Graph}\label{graph}

DET curve: Detection Error Tradeoff curve, 以 FPR （x轴）和 FNR
（y轴）为坐标的曲线。 - DET 越低，模型性能越好。

ROC curve: Receiver Operating Characteristic curve, 以 FPR （x轴）和 TPR
（y轴）为坐标的曲线。 - ROC 越接近左上角，模型性能越好。

AUC: Area Under the Curve, ROC 曲线下的面积，越大越好。

AUC可以看作模型给一个随机正样本打分高于一个随机负样本的概率。（衡量排序能力，AUC
= 1 -\textgreater{} 完美模型）

\begin{itemize}
\tightlist
\item
  Scale-invariant：输出分数范围改变，不影响 AUC
\item
  Threshold-invariant：阈值如何选，不影响 AUC。
\end{itemize}

\[AUC = \frac{1}{m^+ m^-} \sum_{i=1}^{m^+} \sum_{j=1}^{m^-} u(e_{ij})\\
u(e_{ij}) = \begin{cases}
1, & \text{if } s_i^+ > s_j^- \\
0.5, & \text{if } s_i^+ = s_j^- \\
0, & \text{if } s_i^+ < s_j^-
\end{cases}\]


\clearpage

\subsection{Lecture 14: Intro to Unsupervised
Learning}\label{lecture-14-intro-to-unsupervised-learning}

\paragraph{1. Definition}\label{definition}

\begin{longtable}[]{@{}
  >{\raggedright\arraybackslash}p{(\columnwidth - 2\tabcolsep) * \real{0.5000}}
  >{\raggedright\arraybackslash}p{(\columnwidth - 2\tabcolsep) * \real{0.5000}}@{}}
\toprule\noalign{}
\begin{minipage}[b]{\linewidth}\raggedright
Unsupervised Learning
\end{minipage} & \begin{minipage}[b]{\linewidth}\raggedright
Supervised Learning
\end{minipage} \\
\midrule\noalign{}
\endhead
\bottomrule\noalign{}
\endlastfoot
无标签数据 \(D = {(x_i)}_{i=1}^N\) & 有标签数据
\(D = {(x_i, y_i)}_{i=1}^N\) \\
从数据中发现结构、模式、分布等 & 学习输入与输出之间的映射关系 \\
\end{longtable}

\paragraph{2. Main Approaches}\label{main-approaches}

\begin{itemize}
\tightlist
\item
  \textbf{Clustering:}
  把数据点按相似性分成若干组，使组内点相互接近、组间点相互远离。(e.g.~K-means,
  DBSCAN)
\item
  \textbf{Dimensionality Reduction:}
  将高维数据映射到低维空间，同时尽量保留原始数据的结构和信息。(e.g.~PCA,
  t-SNE)
\item
  \textbf{Density Estimation:}
  估计数据的概率分布，常用于异常检测和生成模型。(e.g.~Gaussian Mixture
  Models)
\item
  \textbf{Autoencoders:}
  通过神经网络学习数据的压缩表示，常用于降维和生成任务。
\item
  \textbf{Self-supervised Learning:}
  利用数据本身的结构进行无监督学习，常用于预训练和特征提取。
\end{itemize}


\clearpage

\subsection{Lecture 15: K-means
Clustering}\label{lecture-15-k-means-clustering}

K-means clustering minimizes withincluster variances (squared Euclidean
distances).

本质：把 n 个样本分成 k
组，使得每个样本属于离它最近的那个簇中心所在的组。

\textbf{算法步骤}

\begin{enumerate}
\def\labelenumi{\arabic{enumi}.}
\tightlist
\item
  选定 K，随机初始化 K 个质心（centroids）
\item
  赋值步骤（Assignment）：把每个样本分配给离它最近的质心
\item
  更新步骤（Refitting）：重新计算每个簇的质心（簇内所有点的均值）
\item
  重复2-3直到分配不再改变
\end{enumerate}

\textbf{数学本质} K-means minimize the within-cluster variance：
\[ \begin{aligned}
\min_{c,r} \mathcal{J}(c,r) &= \min_{c,r} \sum_{i=1}^n \sum_{k=1}^K r_{i,k} \| x_i - c_k \|^2 \\
\text{s.t. } r_i &\in \{0,1\}^{n \times K}, \sum_{k=1}^K r_{i,k} = 1
\end{aligned} \] where \(r_{i,k} = 1\) if \(x_i\) is assigned to cluster
\(k\).

这个问题同时优化两类变量（分配 \(r\) 和质心
\(c\)），联合优化很难，于是用 \textbf{坐标下降（Coordinate Descent）}
轮流固定一个优化另一个：

当固定 \(c\) 时，最优的 \(r\) 是将每个样本分配给最近的质心。即
\(k^* = \arg\min_k \| x_i - c_k \|^2\)。

当固定 \(r\) 时，最优的
\(c_k = \frac{\sum_{i=1}^n r_{i,k} x_i}{\sum_{i=1}^n r_{i,k}}\)（簇内样本均值）。

这正是 K-means 算法的赋值步骤和更新步骤。

\textbf{Performance Evaluation of Clustering}

\begin{enumerate}
\def\labelenumi{\arabic{enumi}.}
\item
  Internal Evaluation Metrics.

  只使用数据本身的信息来评估聚类质量。(Silhouette Coefficient).

  对每个样本计算：

  \begin{itemize}
  \tightlist
  \item
    a: 样本与同簇内其他样本的平均距离
  \item
    b: 样本与最近的其他簇内样本的平均距离
  \item
    Silhouette Coefficient: \(s = \frac{b - a}{\max(a, b)}\)，范围
    {[}-1, 1{]}，越接近 1 越好。
  \end{itemize}
\item
  External Evaluation Metrics.

  需要已知真实标签。

  RI 统计所有点对中''分类决策一致''的比例（同簇同类 +
  不同簇不同类），ARI 是对 RI 的随机校正版本，解决了 RI
  对随机聚类也可能较高的问题。

  \begin{itemize}
  \tightlist
  \item
    a： 在 X 中同簇，在 Y 中也同簇
  \item
    b：在 X 中不同簇，在 Y 中也不同簇
  \item
    c：在 X 中同簇，在 Y 中不同簇
  \item
    d：在 X 中不同簇，在 Y 中同簇
  \item
    \(RI = \frac{a + b}{a + b + c + d}\)
    （一致的样本对/总样本对），取值在 {[}0, 1{]}，越接近 1 越好。
  \end{itemize}

  ARI：对 RI 做 chance adjustment，解决了 RI
  对随机聚类也可能较高的问题。
  \[ARI = \frac{RI - E[RI]}{\max(RI) - E[RI]}\] 其中 \(E[RI]\)
  是随机聚类的期望 RI，\(\max(RI)\) 是 RI
  的最大值（当聚类完全正确时）。ARI 的取值范围是 {[}-1, 1{]}，越接近 1
  越好。
\end{enumerate}


\clearpage

\subsection{Lecture 16: Gaussian Mixture Models and
Expectation-Maximization
Algorithm}\label{lecture-16-gaussian-mixture-models-and-expectation-maximization-algorithm}

\paragraph{A. Gaussian Mixture Models
(GMM)}\label{a.-gaussian-mixture-models-gmm}

对于普通分类问题： \$ p(x,z) = p(z) p(x\textbar z) \$ 其中 \(z\)
是类别标签，\(p(z)\) 是类别的先验分布。

但对于聚类问题：z 未知。 对 z 求和：
\[p(x) = \sum_{z} p(x,z) = \sum_{z} p(z) p(x|z)\] 这就是Mixture Model。

\begin{enumerate}
\def\labelenumi{\arabic{enumi}.}
\item
  Gaussian Mixture Model 假设每个簇都是Gaussian分布：
  \[p(x|z=k) = \mathcal{N}(x|\mu_k, \Sigma_k)\] 整体：
  \[p(x) = \sum_{k=1}^K \pi_k \mathcal{N}(x|\mu_k, \Sigma_k)\] 其中
  \(\pi_k\) 是第 k 个簇的权重，满足 \(\sum_{k=1}^K \pi_k = 1\)。
  \[\mathcal{N}(x|\mu, \Sigma) = \frac{1}{(2\pi)^{d/2} |\Sigma|^{1/2}} \exp\left(-\frac{1}{2}(x-\mu)^T \Sigma^{-1} (x-\mu)\right)\]
\item
  GMM 为什么比 K-means 更强大？
\end{enumerate}

\begin{itemize}
\tightlist
\item
  K-means 只能捕捉球状簇，而 GMM 可以捕捉任意形状的簇。
\item
  GMM 提供了每个样本属于每个簇的概率，而 K-means 只能给出硬分配。
\end{itemize}

\begin{enumerate}
\def\labelenumi{\arabic{enumi}.}
\setcounter{enumi}{2}
\tightlist
\item
  \textbf{参数估计} 目标：求
  \(\theta = \{\pi_k, \mu_k, \Sigma_k\}_{k=1}^K\)
  使得数据的似然函数最大化：
  \[\mathcal{L} = \log L(\Theta)=\sum_{n=1}^{N}\log \left(\sum_{k=1}^{K}\pi_k \mathcal{N}(x^{(n)}|\mu_k,\Sigma_k)\right)\]
  where \(\Theta = \{\boldsymbol{\pi}, \boldsymbol{\mu}, \boldsymbol{\Sigma}\}\).
\end{enumerate}

\[
\frac{\partial \mathcal{L}}{\partial \boldsymbol{\mu}_k} = 0
\implies
\sum_{n=1}^N
\underbrace{
\frac{
\pi_k \mathcal{N}(\mathbf{x}^{(n)} \mid \boldsymbol{\mu}_k, \boldsymbol{\Sigma}_k)
}{
\sum_j \pi_j \mathcal{N}(\mathbf{x}^{(n)} \mid \boldsymbol{\mu}_j, \boldsymbol{\Sigma}_j)
}
}_{\gamma_k^{(n)}}
\boldsymbol{\Sigma}_k^{-1} (\mathbf{x}^{(n)} - \boldsymbol{\mu}_k)
= 0
\]

引入 responsibility 责任度 \(\gamma_k^{(n)}\)：
\[\gamma_k^{(n)} = \frac{\pi_k \mathcal{N}(x^{(n)}|\mu_k, \Sigma_k)}{\sum_{j=1}^K \pi_j \mathcal{N}(x^{(n)}|\mu_j, \Sigma_j)}\]

\(\gamma_k^{(n)}\) 表示样本 \(x^{(n)}\) 属于簇 \(k\) 的后验概率。

\begin{enumerate}
\def\labelenumi{\arabic{enumi}.}
\setcounter{enumi}{3}
\tightlist
\item
  \textbf{参数更新共识} 有了责任度后，可以得出以下更新公式：
\end{enumerate}

\begin{itemize}
\tightlist
\item
  有效样本数：\(N_k = \sum_{n=1}^N \gamma_k^{(n)}\)
\item
  均值更新：\(\mu_k = \frac{1}{N_k} \sum_{n=1}^N \gamma_k^{(n)} x^{(n)}\)
\item
  协方差更新：\(\Sigma_k = \frac{1}{N_k} \sum_{n=1}^N \gamma_k^{(n)} (x^{(n)} - \mu_k)(x^{(n)} - \mu_k)^T\)
\item
  混合系数更新：\(\pi_k = \frac{N_k}{N}\)
\end{itemize}

\begin{enumerate}
\def\labelenumi{\arabic{enumi}.}
\setcounter{enumi}{4}
\item
  \textbf{交替更新算法} 第一步：计算责任度
  \[\gamma_k^{(n)} = \frac{\pi_k \mathcal{N}(x^{(n)}|\mu_k, \Sigma_k)}{\sum_{j=1}^K \pi_j \mathcal{N}(x^{(n)}|\mu_j, \Sigma_j)}\]
  第二步：更新参数 \[N_k = \sum_{n=1}^N \gamma_k^{(n)} \\
  \mu_k = \frac{1}{N_k} \sum_{n=1}^N \gamma_k^{(n)} x^{(n)} \\
  \Sigma_k = \frac{1}{N_k} \sum_{n=1}^N \gamma_k^{(n)} (x^{(n)} - \mu_k)(x^{(n)} - \mu_k)^T \\
  \pi_k = \frac{N_k}{N}\] 不断重复直到对数似然收敛。
\item
  Comparison with K-means
\end{enumerate}

\begin{itemize}
\tightlist
\item
  K- Means

  \begin{itemize}
  \tightlist
  \item
    Assignment step：把每个样本分给最近的簇中心（硬分配）
  \item
    Refitting step：更新簇中心为簇内样本的均值
  \end{itemize}
\item
  GMM

  \begin{itemize}
  \tightlist
  \item
    E-step：计算每个样本属于每个簇的后验概率（软分配）
  \item
    M-step：根据责任度更新簇的参数（均值、协方差、混合系数）
  \end{itemize}
\end{itemize}

\paragraph{B. Expectation-Maximization (EM)
Algorithm}\label{b.-expectation-maximization-em-algorithm}

\begin{enumerate}
\def\labelenumi{\arabic{enumi}.}
\tightlist
\item
  GMM 的隐变量视角 We introduce a latent variable \(z\) for each sample,
  where \(z_i\) indicates which component generated \(x_i\).
\end{enumerate}

\[z \sim \text{Categorical}(\pi) \\
p(z = k) = \pi_k \\
p(x|z=k) = \mathcal{N}(x|\mu_k, \Sigma_k)\]

于是：
\[p(x) = \sum_{k=1}^K p(z=k) p(x|z=k) = \sum_{k=1}^K \pi_k \mathcal{N}(x|\mu_k, \Sigma_k)\]
这样就把 GMM 解释成一个含离散隐变量的生成模型。

\(z^{(n)}\) 不知道，难以用 MLE 直接求解参数。

\begin{enumerate}
\def\labelenumi{\arabic{enumi}.}
\setcounter{enumi}{1}
\tightlist
\item
  Jensen's Inequality Suppose \(f\) is a convex function, and X is a
  random variable, then: \[f(\mathbb{E}[X]) \leq \mathbb{E}[f(X)]\] If
  \(f\) is concave, the inequality is reversed:
  \[f(\mathbb{E}[X]) \geq \mathbb{E}[f(X)]\]
\end{enumerate}

\begin{itemize}
\tightlist
\item
  E-step：先猜每个样本来自各个簇的概率
\item
  M-step：假设这些概率是真的，再重新估计参数
\end{itemize}

\end{document}
